\section{Siglas de interés}
\begin{itemize}
    \item \textbf{ASIC}: \textit{Application-Specific Integrated Circuit}, son circuitos integrados diseñados para una función específica y rígida, no programables. \cite{harris}
    \item \textbf{PAL}: \textit{Programmable Array Logic}: Dispositivos programables mediante la quema de fusibles (antecesor de las FPGA).
    \item \textbf{PLL/MMCM}: Circuitos que actúan como ''directores de orquesta'', generando señales de reloj precisas.
    \item \textbf{SoC}: \textit{System on a Chip}, chips que integran múltiples componentes en una sola pieza de silicio.
    \item \textbf{RISC-V}: Arquitectura de conjunto de instrucciones de código abierto muy usada actualmente dentro de los FPGA.
    \item \textbf{AXI}: \textit{Advanced eXtensible Interface} bus estándar de comunicación interna de alta velocidad.
    \item \textbf{SSI}: \textit{Stracked Silicon Interconnect} Tecnología para apilar trozos de silicio, precursora de modernos chipets.
    \item \textbf{DRP}: \textit{Dynamic Partial Reconfiguration} Capacidad de cambiar una parte de la lógica del chip mientras el resto sigue operando.
    \item \textbf{HFT}: \textit{High Frequency Trading} algoritmos de trading de alta frecuencia que utilizan FPGA (ámbito financiero).
    \item \textbf{ACAP}: \textit{Adaptive Compute Acceleration Platform}, evolución avanzada de los FPGA que integra motores de inteligencia artifical dedicados.
\end{itemize}
